\setcounter{page}{137}
\chapter{DUALITY FOR PROJECTIVE MORPHISMS}
\section{\S1 Differentials}

This section reviews basic facts about relative differentials; standard references are SGA 60--61 Exposé II (cf.\ future EGA treatments).

\textbf{Definition.}
Let $A\to B$ be a ring homomorphism, $M$ a $B$-module.
Define $\Der_A(B,M)$ as the $A$-module of $A$-derivations $B\to M$.
The \textit{module of relative 1-differentials} $\Omega^1_{B/A}$ is the $B$-module representing the functor
\[
M\mapsto \Der_A(B,M).
\]
Equivalently, there is a universal derivation $d\colon B\to\Omega^1_{B/A}$ such that the natural map
\[
\Hom_B(\Omega^1_{B/A},M) \to \Der_A(B,M)
\]
is an isomorphism for all $B$-modules $M$.

For a morphism of preschemes $f\colon X\to Y$, the sheaf of relative 1-differentials $\Omega^1_{X/Y}$ is defined by gluing the affine local data $\Omega^1_{B/A}$ over affine open coverings of $X$ and $Y$.

\setcounter{page}{138}
\textbf{Definition} [EGA IV 6.8.1].
A morphism $f\colon X\to Y$ is \textit{smooth} if it is flat, locally of finite presentation, and every fibre $f^{-1}(y)$ is locally noetherian and geometrically regular.

\textbf{Examples.}
1. Open immersions are smooth.
2. Compositions of smooth morphisms are smooth.
3. Smooth morphisms are stable under base change.
4. A prescheme over a field is smooth iff it is locally noetherian and geometrically regular.

\textbf{Proposition 1.1} [SGA 60--61, II.4.3].
Let $f\colon X\to Y$, $g\colon Y\to S$ be smooth morphisms of preschemes.
There is a canonical exact sequence
\[
0 \to f^*\Omega^1_{Y/S} \to \Omega^1_{X/S} \to \Omega^1_{X/Y} \to 0,
\]
and $\Omega^1_{X/Y}$ is locally free of rank equal to the relative dimension of $X$ over $Y$.

\textbf{Definition.}
A closed subscheme $Y\subseteq X$ is \textit{locally a complete intersection} if every point $y\in Y$ has an open neighbourhood $U\subseteq X$ such that the ideal sheaf $\mathcal{J}_Y|_U$ is generated by an $\sheaf{O}_X$-regular sequence $s_1,\dots,s_r$:
$s_1$ is a non-zero divisor, and for each $2\le i\le r$, $s_i$ is a non-zero divisor modulo $(s_1,\dots,s_{i-1})$.

\setcounter{page}{139}
\textbf{Proposition 1.2} [SGA 60--61, II.10].
Let $X$ be locally noetherian and smooth over $S$, and let $Y\subseteq X$ be closed.
The following are equivalent:
\begin{enumerate}
\item[(i)] $Y$ is smooth over $S$;
\item[(ii)] $\Omega^1_{Y/S}$ is locally free, and the sequence
\[
0 \to \mathcal{J}/\mathcal{J}^2 \to i^*\Omega^1_{X/S} \to \Omega^1_{Y/S} \to 0
\]
is exact, where $i\colon Y\hookrightarrow X$ is the closed immersion and $\mathcal{J}$ the ideal sheaf of $Y$.
\end{enumerate}
In particular, such $Y$ is locally a complete intersection in $X$.

\textbf{Definition.}
Given an exact sequence of locally free sheaves
\[
(*) \quad 0\to E'\to E\to E''\to 0
\]
of ranks $r,\,r+s,\,s$ respectively, define a canonical isomorphism
\[
\varphi(*)\colon \Lambda^{r+s}E \xrightarrow{\sim} \Lambda^r E' \otimes \Lambda^s E''.
\]
Choose a local basis $e_1,\dots,e_{r+s}$ of $E$ such that $e_1,\dots,e_r$ trivialize $E'$ and the images of $e_{r+1},\dots,e_{r+s}$ trivialize $E''$. The isomorphism is
\[
e_1\wedge\cdots\wedge e_{r+s}
\mapsto
(e_1\wedge\cdots\wedge e_r)\otimes (\overline{e}_{r+1}\wedge\cdots\wedge\overline{e}_{r+s}).
\]

\setcounter{page}{140}
\textbf{Remark.}
These isomorphisms depend on choices; compatibilities are not automatic. For dual sequences, the induced isomorphism differs by a sign $(-1)^{rs}$.

\textbf{Lemma 1.3.}
Let $0\subseteq E_1\subseteq E_2\subseteq E$ be a flag of locally free sheaves on $X$.
The four short exact sequences associated to the flag induce a commutative diagram of the exterior power isomorphisms $\varphi(*)$.

\textit{Proof.} Omitted.

\setcounter{page}{141}
\textbf{Definition.}
\begin{enumerate}
\item[(a)] Let $f\colon X\to Y$ be smooth of relative dimension $n$. Define
\[
\omega_{X/Y} = \Lambda^n \Omega^1_{X/Y}.
\]
By Proposition 1.1, $\omega_{X/Y}$ is invertible (rank~1) locally free.

\item[(b)] Let $f\colon X\hookrightarrow Y$ be a closed immersion locally a complete intersection of codimension $n$. Let $\mathcal{J}$ be the ideal sheaf of $X$ in $Y$. Define
\[
\omega_{X/Y} = \big(\Lambda^n (\mathcal{J}/\mathcal{J}^2)\big)^\vee,
\]
where $(-)^\vee$ denotes dual sheaf. Since $\mathcal{J}/\mathcal{J}^2$ is locally free of rank $n$, $\omega_{X/Y}$ is invertible.
\end{enumerate}

\textbf{Remarks.}
1. A morphism that is both smooth and a closed immersion is locally an isomorphism, so the two definitions agree.
2. If $X'\to Y'$ is a base change of $X\to Y$, then $\mathrm{pr}_1^*\omega_{X/Y} \cong \omega_{X'/Y'}$.

\setcounter{page}{142}
\textbf{Lemma 1.4.}
Let $X\xrightarrow{f}Y\xrightarrow{g}Z$ be prescheme morphisms with $g$ smooth.
Let $\Gamma\colon X\to X\times_Z Y$ be the graph morphism. Then
\[
\omega_{X/X\times_Z Y} \cong f^*\omega_{Y/Z}.
\]
\textit{Proof.}
The graph $\Gamma$ is a local complete intersection. Apply base change and the canonical isomorphism for complete intersection dualizing sheaves.

\setcounter{page}{143}
\textbf{Definition 1.5.}
Let $X\xrightarrow{f}Y\xrightarrow{g}Z$ be composable morphisms of locally noetherian preschemes, each either smooth or locally a complete intersection. There is a canonical isomorphism
\[
\zeta_{f,g}\colon \omega_{X/Z} \xrightarrow{\sim} f^*\omega_{Y/Z} \otimes \omega_{X/Y}.
\]
There are four cases:
\begin{enumerate}
\item[(a)] $f,g,gf$ all smooth: use the exterior power isomorphism on the relative differential exact sequence.
\item[(b)] $f,g,gf$ all local complete intersections: use the ideal sheaf exact sequence and dualize exterior powers.
\item[(c)] $f$ l.c.i., $g,gf$ smooth.
\item[(d)] $f,gf$ l.c.i., $g$ smooth: use the graph construction and Lemma 1.4.
\end{enumerate}

\textbf{Proposition 1.6.}
For three composable morphisms $X\xrightarrow{f}Y\xrightarrow{g}Z\xrightarrow{h}W$ (each smooth or l.c.i.), the diagram
\[
\zeta_{h,g}\circ\zeta_{f,hg} = \zeta_{f,g}\circ\zeta_{gf,h}
\]
commutes.

\textit{Proof.} Follows from Lemma 1.3.

\setcounter{page}{144}
\textbf{Remark.}
The proper framework for $\omega_{X/Y}$ is the theory of \textit{Gorenstein morphisms}:
A locally finite-type morphism $f\colon X\to Y$ of locally noetherian preschemes is Gorenstein if it has finite Tor-dimension and $Rf^*\sheaf{O}_Y\in D^+(X)$ is invertible. This invertible sheaf is defined to be $\omega_{X/Y}$. Smooth and l.c.i.\ morphisms are Gorenstein. For composable Gorenstein morphisms the same multiplicative isomorphism $\zeta_{f,g}$ exists, and the triple composition diagram commutes.